\addcontentsline{toc}{chapter}{Abstract}
\section*{\centering Abstract}

Historical musical instruments present a complex set of challenges within the
heritage sector, as they span across many categories of value systems in
heritage conservation. The research in this thesis was carried out as part of
the NEMUS project, which seeks to address such challenges presented by musical
instrument heritage objects, through the application of digital technologies. Two
challenges of digital musical instrument heritage conservation are explored:
interaction with digitised instruments and analysis for understanding material
properties. 

The challenge of interaction with digitised instruments is explored through the
harpsichord, for which standard MIDI-based controllers fail to capture the
distinct haptic experience during performance. This challenge is explored
through the creation a two-register haptic interface constructed from a
custom-built, historically informed harpsichord keyboard the provides the same
kinaesthetic response of an original instrument. The interface's design
principles, prototyping and fabrication process are discussed alongside the open
methodology which informed them, culminating in an exhibition at Museo San
Colombano, Bologna. 

To address the challenge of analysing material properties, an open-source finite
difference simulation framework for plate vibrations featuring adjustable
elastic boundary conditions was created. The design and development of the
software, titled \textsc{magpie}, are discussed with particular attention paid
to its cross-language implementation in M\textsc{atlab}, Python and
C\textsc{++}. \textsc{magpie} aims to democratizes access to acoustic modelling,
benefiting researchers, educators, and instrument builders by visualizing mode
shapes and frequencies, facilitating experimental exploration of the plate
systems associated with string instrument soundboards. 

Inherent to both projects is an open methodology, an adherence to FAIR
principles in digital research. Digital technologies can support the
conservation, analysis and recreation of historical musical instruments. With
this technology, comes the challenges of sustaining digital artefacts,
reproducibility, software citation and sustainable data practices. The future
directions for each project are examined, with prospective developments and
technical challenges contextualised with the broader open methodology.


\cleardoublepage{}